\chapter{文献严格化：$t>3$ 门槛与剔除微盘}\label{app:strict}

\section{为什么单开一章}

2026 年 11 月，一份外部讲义（\textit{AI 炒股与量化交易：成熟技术讲义}，
16 章 + 4 附录）提供了两条对本项目有直接冲击的文献结论：

\begin{enumerate}
  \item \textbf{Harvey, Liu \& Zhu (2016)}：在数百个已发表因子之后，
        $t>2.0$ 的显著性门槛已不够，新因子应要求 $t>3.0$。
  \item \textbf{Liu, Stambaugh \& Yuan (2019)}：A 股的小市值效应被"壳价值"污染，
        建议剔除最小 30\% 市值后再检验。
\end{enumerate}

本书此前的所有因子判定都用 $t>2$，也从未做过剔除微盘的稳健性检验。
这一章把两条建议变成可执行的检验。

\begin{warnbox}[样本口径说明]
本章的两张表用了不同区间：严格性检验（表~\ref{tab:strict}）用 2013--2026，
因为需要 20 日滚动变量；$t$ 门槛重判用全样本 2010--2026（来自第~\ref{ch:result} 章的因子表）。
同一个因子在两表里的数值略有差异，源于区间不同，不是错误。
\end{warnbox}

\section{检验一：剔除最小 30\% 市值}

按当日对数流通市值分位，剔除最小的 30\%，在剩余样本上重算 Rank IC。

\begin{table}[htbp]
\centering
\caption{全样本 vs 剔除最小 30\% 市值（$h=5$，2013--2026）}
\label{tab:strict}
\small
\begin{tabular}{@{}lrrrrrr@{}}
\toprule
因子 & IC 全样本 & IC 剔除小30\% & IC 变化 & $t$ 全样本 & $t$ 剔除后 & 判定 \\
\midrule
max\_ret\_20 & +0.0592 & +0.0538 & -9.1\% & +17.29 & +13.94 & 稳健 \\
turn\_20 & +0.0585 & +0.0504 & -13.8\% & +15.44 & +12.61 & 稳健 \\
pv\_corr\_20 & +0.0346 & +0.0310 & -10.5\% & +15.03 & +12.70 & 稳健 \\
skew\_20 & +0.0260 & +0.0202 & -22.2\% & +13.44 & +10.17 & 减弱 \\
vol\_20 & +0.0606 & +0.0566 & -6.6\% & +12.81 & +10.80 & 稳健 \\
ln\_amount\_20 & +0.0590 & +0.0551 & -6.6\% & +12.68 & +10.43 & 稳健 \\
mom\_20 & -0.0500 & -0.0437 & -12.6\% & -12.22 & -10.08 & 稳健 \\
ivol\_60 & +0.0698 & +0.0677 & -3.0\% & +11.10 & +9.72 & 稳健 \\
rev\_10 & +0.0380 & +0.0323 & -15.1\% & +10.43 & +8.47 & 稳健 \\
rev\_5 & +0.0308 & +0.0259 & -16.1\% & +9.63 & +7.69 & 稳健 \\
rev\_1 & +0.0159 & +0.0135 & -15.2\% & +8.66 & +7.00 & 稳健 \\
bp & +0.0339 & +0.0317 & -6.6\% & +8.50 & +6.55 & 稳健 \\
amihud\_20 & +0.0231 & +0.0149 & -35.6\% & +7.46 & +3.88 & 减弱 \\
mom\_120 & -0.0412 & -0.0389 & -5.7\% & -6.24 & -5.26 & 稳健 \\
size & +0.0097 & -0.0050 & -152.2\% & +1.26 & -0.76 & \textbf{翻负，不可用} \\
\bottomrule
\end{tabular}

\end{table}

\subsection{怎么读}

\begin{enumerate}
  \item \textbf{反转、波动、换手、流动性四族非常稳健}。\texttt{ivol\_60} 只掉 $3.0\%$，
        \texttt{vol\_20} 与 \texttt{ln\_amount\_20} 各掉 $6.6\%$，
        剔除微盘之后 $t$ 值仍在 $9.7$--$10.8$。这说明它们的收益不依赖最小的那三成。
  \item \textbf{短期反转减弱但存活}。\texttt{rev\_5} 掉 $16.1\%$，$t$ 从 $9.63$ 降到 $7.69$。
        微盘确实贡献了一部分反转收益，但远不是全部。
  \item \textbf{Amihud 非流动性掉得最多}（$-35.6\%$）。这符合逻辑：
        非流动性溢价本来就集中在流动性最差的股票上，而它们正是被剔除的那批。
  \item \textbf{规模因子直接翻负}。\texttt{size} 的 IC 从 $+0.0097$ 变成 $-0.0050$，
        $t$ 从 $+1.26$ 变成 $-0.76$。
\end{enumerate}

\begin{keybox}[这一条最重要]
\texttt{size} 的正 IC \textbf{100\% 来自最小那 30\% 市值}。
一旦剔除，它不仅消失，还翻成负号。

这意味着：本项目此前把规模因子列为"衰减"是过于客气的说法。
准确的说法是——\textbf{它在可执行的那部分市场里根本不存在}。

它与三件事同时吻合：LSY (2019) 的"壳价值污染"、外部讲义第 13 章的
"2024 年初微盘股流动性危机"，以及附录~\ref{app:shm} 发现的
"外部快照恰好缺最小那三成"。三者指向同一个结论：
\textbf{微盘股上的"规律"既不可执行，也不可信。}
\end{keybox}

\section{检验二：把门槛提到 $t>3$}

按 Harvey-Liu-Zhu 的建议，用 $|t|>3$ 重新判定第~\ref{ch:result} 章的 39 个因子：

\begin{tabular}{@{}lr@{}}
\toprule
判定 & 因子数 \\
\midrule
$|t|>3$（HLZ 建议门槛） & \textbf{31} \\
仅 $2<|t|\le 3$ & 4 \\
$|t|\le 2$ & 4 \\
\bottomrule
\end{tabular}


仅通过 $t>2$、未过 $t>3$ 的四个因子：

\begin{tabular}{@{}llrr@{}}
\toprule
因子 & 家族 & IC & $t$ \\
\midrule
rev\_yoy & 成长 & +0.0084 & +2.94 \\
size & 规模 & +0.0167 & +2.27 \\
kurt\_20 & 波动 & +0.0050 & +2.13 \\
mom\_12\_1 & 动量 & -0.0687 & -2.07 \\
\bottomrule
\end{tabular}


\begin{keybox}[结论]
39 个因子里，\textbf{31 个通过 $|t|>3$}。这个比例远高于 Hou-Xue-Zhang (2020)
对已发表异象的复现结果（他们发现 447 个异象中 64\% 在 5\% 水平下不显著）。

两个可能的原因：
\begin{enumerate}
  \item 本项目的因子集中在\textbf{量价类}（反转、波动、换手、流动性），
        这一类在 A 股本来就更强，而 HLZ/HXZ 的样本以美股与基本面因子为主；
  \item $t$ 值是在\textbf{全序列}上算的 HAC $t$（观测数 3400 左右），
        而不是"各年 $t$ 的平均"，前者天然更大。
\end{enumerate}

因此这 31 个不应被理解为"31 个独立的新发现"，
而应理解为"31 个因子在样本内具有统计上难以用噪声解释的排序能力"。
它们之间的相关性很高（第~\ref{ch:factor} 章已讨论），
真正独立的维度远小于 31。
\end{keybox}

\section{对被淘汰因子与规模因子的处理}

\begin{enumerate}
  \item \textbf{规模因子（\texttt{size}）降级为"不可用"}：
        它只在最小 30\% 市值里有效，而那一区段容量极低、2024 年出现过流动性危机，
        且外部数据快照根本不覆盖。因子库保留其代码以备复现，但不再进入任何合成。
  \item \textbf{Amihud 与偏度标记为"减弱"}：剔除微盘后仍显著，但幅度掉三成以上，
        使用时必须注明"含微盘暴露"。
  \item \textbf{四个仅过 $t>2$ 的因子}（\texttt{rev\_yoy}、\texttt{size}、
        \texttt{kurt\_20}、\texttt{mom\_12\_1}）在结论中一律标注为"边缘"，
        不作为核心证据引用。
\end{enumerate}

\section{附：证据薄弱清单与识别红旗}

以下内容改编自外部讲义第 14 章，用于快速排除不合格的说法。

\begin{center}
\small
\begin{tabular}{@{}p{0.30\textwidth}p{0.60\textwidth}@{}}
\toprule
流行说法 & 证据状况与替代方案 \\
\midrule
"深度学习能预测股价" & 证据薄弱：非平稳 + 扣成本后优势消失。
改用 ML 做横截面排序并控制换手。 \\
"LLM 选股年化 30\%+" & 预训练记忆污染 + 不可交易。
只用模型 cutoff 之后的样本；做事件抽取而不是选股。 \\
"强化学习自动交易" & 样本效率与奖励黑客问题。用执行与做市，不用于选股。 \\
"因子越多越强" & 已被 HLZ (2016) 与 HXZ (2020) 证伪。少数稳健因子 + 风险模型。 \\
"夏普 3 的策略" & 大概率没算 DSR/PBO。先问"试了多少次"。 \\
"十年回测稳定盈利" & 多半未含退市股。要求全样本 + PIT。 \\
"高频稳赚" & 军备竞赛 + 监管 + 成本。转向中低频与成本控制。 \\
\bottomrule
\end{tabular}
\end{center}

\begin{warnbox}[识别"AI 炒股"产品的六个红旗]
\begin{enumerate}
  \item 只展示收益曲线，不展示回撤、换手率、容量；
  \item 拿不出连续 3 年以上的\textbf{实盘}（非模拟）记录；
  \item 拒绝说明策略逻辑与风险来源，只讲"神经网络""大模型"；
  \item 承诺保本或固定收益（在中国属于违规宣传）；
  \item 收益曲线过于平滑（夏普 $>4$ 且回撤 $<3\%$），真实市场几乎不存在；
  \item 靠卖软件、卖课、卖会员盈利，而不是靠自营交易盈利。
\end{enumerate}

一个可执行的检验：让对方用"未来 6 个月的纸上交易"证明，而不是展示过去。
\end{warnbox}

\section{本章小结}

\begin{enumerate}
  \item 剔除最小 30\% 市值后，反转、波动、换手、流动性四族仍然高度显著，
        IC 只下降 3\%--16\%。
  \item 规模因子的正 IC 全部来自最小那三成，剔除后翻负——
        应判定为\textbf{不可用}，而不是"衰减"。
  \item Amihud 非流动性下降 36\%，属于"减弱"，使用时须注明微盘暴露。
  \item 按 $|t|>3$ 重判，39 个因子中 31 个通过，4 个仅在 $t>2$ 水平上显著。
  \item "31 个通过 $t>3$"不等于 31 个独立发现；因子之间高度相关，
        真正独立的维度要少得多。
\end{enumerate}
